\chapter{LLM이란 무엇인가}
\label{ch:intro}

이 장에서는 대형 언어 모델(Large Language Model, LLM)이 무엇인지, 어떻게 발전해 왔는지, 그리고 왜 직접 만들어보는 것이 가치 있는지를 살펴본다. 코드보다 개념에 집중하므로, 딥러닝을 처음 접하는 독자도 부담 없이 읽을 수 있다.

\section{언어 모델이란}

\keyterm{언어 모델(language model)}은 단어(또는 토큰)의 시퀀스에 확률을 할당하는 통계적 모델이다. 쉽게 말해, ``이 문장이 자연스러운가?''를 확률로 판단하는 함수 $P(w_1, w_2, \dots, w_n)$다. 예를 들어 ``고양이가 마우스를 쫓는다''는 자연스럽지만 ``마우스가 고양이를 쫓는다''는 문맥에 따라 덜 자연스러울 수 있다. 언어 모델은 이런 확률을 학습한다.

역사적으로 언어 모델은 N-gram 모델에서 시작해 신경망 기반 모델(RNN, LSTM)을 거쳐, 2017년 트랜스포머(Transformer)의 등장과 함께 비약적으로 발전했다. 특히 ``다음 토큰 예측(next token prediction)''이라는 단순한 목표로 인터넷 규모의 텍스트를 학습한 모델들이 놀라운 언어 이해 능력을 보여주면서, 오늘날의 LLM 시대가 열렸다.

\begin{infobox}[다음 토큰 예측]
LLM의 학습 목표는 단순하다: 지금까지의 토큰들을 보고 다음 토큰을 예측하는 것. 수식으로 쓰면
\[
P(w_t \mid w_1, w_2, \dots, w_{t-1})
\]
이 조건부 확률을 정확히 추정하는 것이 언어 모델의 유일한 목표다. 이 단순한 목표가 대규모 데이터와 결합했을 때 질문 응답, 코드 작성, 번역 등 복잡한 능력이 자연스럽게 등장한다.

왜 이렇게 단순한 목표가 강력한가? 언어를 ``이해''하려면 문법, 상식, 추론, 사실 지식 등을 모두 갖춰야 한다. 다음 토큰을 정확히 예측하려면 모델은 이 모든 것을 암묵적으로 학습해야 한다. ``지구는 태양 주위를 도는 ...'' 다음에 ``행성''을 예측하려면 천문학 상식이 필요하고, ``1+1='' 다음에 ``2''를 예측하려면 산술이 필요하다.
\end{infobox}

\section{직관: 언어 모델을 어떻게 생각할 것인가}

언어 모델을 처음 접하는 독자를 위해 비유를 들어보자. 언어 모델은 ``엄청나게 많은 책을 읽은 사람''과 비슷하다. 이 사람은 모든 문장의 ``다음 단어''를 맞추는 훈련을 받았다. 훈련이 끝나면 그 사람은 질문을 받았을 때 가장 자연스러운 다음 단어들을 이어붙여 답변을 ``생성''한다.

물론 진짜 사람과는 다르다. 언어 모델은 ``이해''한다기보다 ``패턴을 매칭''한다. 하지만 충분히 많은 패턴을 학습하면 겉보기에 이해하는 것처럼 보인다. 이것이 바로 GPT-4 같은 모델이 놀라운 답변을 내놓는 원리다.

\begin{notebox}[언어 모델은 ``이해''하는가?]
이 질문은 철학적 논쟁거리다. 기능주의(functionalism) 관점에서 ``입력에 대해 적절히 반응하면 이해한다''고 볼 수 있고, 반대로 ``내부 표상이 없으면 이해가 아니다''라고 볼 수도 있다. 이 책에서는 철학적 논쟁보다는 ``어떻게 동작하는가''에 집중한다. 철학에 관심이 있다면 Searle의 ``중국어 방'' 사고실험을 찾아보라.
\end{notebox}

\section{LLM의 발전 역사}

\subsection{통계적 언어 모델 (1990년대$\sim$2000년대)}

초기 언어 모델은 N-gram이라는 단순한 통계 기법을 사용했다. ``직전 N-1개의 단어를 보고 다음 단어의 확률을 추정''하는 방식이다. 예를 들어 ``the cat sat on the'' 다음에 올 단어로 ``mat''이 ``banana''보다 확률이 높다는 것을 말뭉치 통계로 학습한다.

\begin{lstlisting}[language=Python]
# N-gram 언어 모델의 핵심 아이디어 (단순화)
from collections import Counter

corpus = "the cat sat on the mat the cat ran".split()
# 2-gram (bigram) 카운트
bigrams = [(corpus[i], corpus[i+1]) for i in range(len(corpus)-1)]
counts = Counter(bigrams)
# P("mat" | "the") 추정
p_mat_given_the = counts[("the", "mat")] / sum(1 for w1, w2 in bigrams if w1 == "the")
print(f"P(mat|the) = {p_mat_given_the:.2f}")
\end{lstlisting}

하지만 N-gram은 심각한 한계가 있었다. 문맥을 짧게만 볼 수 있고(N이 보통 3$\sim$5), 학습하지 않은 조합에는 0의 확률을 할당하는 ``희소성(sparsity)'' 문제가 있었다. 이를 보완하기 위해 smoothing 기법이 개발되었지만 근본적 해결은 아니었다.

또 다른 문제는 ``단어의 의미''를 포착하지 못한다는 것이었다. ``개''와 ``강아지''는 비슷한 의미지만 N-gram에서는 완전히 다른 단어로 취급된다. 이를 해결하기 위해 단어를 분산 표현(distributed representation)으로 바꾸는 신경망 기반 모델이 등장했다.

\subsection{신경망 언어 모델 (2003$\sim$2017)}

Bengio et al. (2003)은 최초로 신경망 기반 언어 모델을 제안했다. 단어를 저차원 밀집 벡터(임베딩)로 표현하고, 이 임베딩들을 신경망에 넣어 다음 단어 확률을 예측하는 방식이었다. 이후 RNN(Recurrent Neural Network)과 LSTM(Long Short-Term Memory)이 등장하면서 더 긴 문맥을 처리할 수 있게 되었다.

\begin{infobox}[Word2Vec과 임베딩]
Mikolov et al. (2013)의 Word2Vec은 단어를 100$\sim$300차원 벡터로 표현하면서 ``비슷한 의미의 단어는 벡터 공간에서 가깝다''는 것을 보였다. 유명한 예: $\vec{\text{king}} - \vec{\text{man}} + \vec{\text{woman}} \approx \vec{\text{queen}}$. 이 아이디어는 4장에서 구현할 임베딩의 기반이다.
\end{infobox}

\begin{warnbox}[RNN의 한계]
RNN은 이전 시점의 은닉 상태를 현재 시점에 전달하는 구조다. 이론적으로는 무한한 문맥을 처리할 수 있지만, 실제로는 역전파 시 그래디언트가 사라지거나 폭발하는 \keyterm{vanishing/exploding gradient} 문제가 발생한다. 이로 인해 실제 처리 가능한 문맥 길이는 수십$\sim$수백 토큰에 불과했다. 또한 순차 처리만 가능해 GPU 병렬화에 비효율적이었다.

LSTM(Long Short-Term Memory)은 게이트 메커니즘으로 이 문제를 부분적으로 완화했지만, 근본적 해결은 아니었다. 100단어 이상의 문맥은 여전히 어려웠다.
\end{warnbox}

\subsection{트랜스포머의 등장 (2017)}

2017년 Vaswani et al.이 발표한 ``Attention Is All You Need''는 언어 모델의 판도를 완전히 바꿨다. RNN 없이 어텐션(attention) 메커니즘만으로 시퀀스를 처리하는 트랜스포머(Transformer) 구조를 제안했다. 트랜스포머의 핵심 혁신은 두 가지다.

\begin{enumerate}
  \item \textbf{병렬 처리}: RNN처럼 순차적으로 처리하지 않고, 시퀀스의 모든 위치를 동시에 처리할 수 있어 GPU 활용도가 극대화되었다. 이는 학습 속도를 수십 배 향상시켰다.
  \item \textbf{장거리 의존성}: 어텐션은 시퀀스의 모든 위치 쌍 사이의 직접 연결을 계산하므로, 거리가 먼 토큰 간의 관계도 잘 포착한다. 1000단어 떨어진 단어라도 한 번의 어텐션으로 참조 가능하다.
\end{enumerate}

\begin{figure}[htbp]
\centering
\begin{tikzpicture}[
  scale=0.9,
  every node/.style={font=\small\sffamily},
  bar/.style={inkblue, fill=inkblue!30},
  lbl/.style={inkgray, anchor=south, font=\scriptsize\sffamily}
]
  % 축
  \draw[inkblue, line width=0.5pt, ->] (0, 0) -- (11, 0);
  \draw[inkblue, line width=0.5pt, ->] (0, 0) -- (0, 5);
  \node[inkgray, anchor=east] at (-0.2, 4.7) {params};
  \node[inkgray, anchor=north] at (10.5, -0.2) {year};
  % 로그 스케일 막대 (대략적)
  \foreach \x/\h/\name in {1/0.2/GPT-1, 2.5/1.2/BERT, 4/1.8/GPT-2, 5.5/3.0/T5, 7/4.0/GPT-3, 8.5/4.5/PaLM, 10/4.8/GPT-4} {
    \draw[bar] (\x-0.3, 0) rectangle (\x+0.3, \h);
    \node[lbl] at (\x, \h+0.05) {};
    \node[inkgray, anchor=south, font=\tiny\sffamily, rotate=45] at (\x, 0.05) {\name};
  }
\end{tikzpicture}
\caption{LLM 파라미터 수의 발전 추이 (로그 스케일, 대략적). GPT-1 (2018)은 1.1억 개, GPT-4 (2023)는 추정 1.7조 개.}
\label{fig:llm-scaling}
\end{figure}

\subsection{GPT 시리즈와 스케일링 법칙 (2018$\sim$현재)}

OpenAI의 GPT 시리즈는 트랜스포머 디코더 구조를 사용해 대규모 텍스트로 사전학습(pre-training)하는 접근을 취했다. GPT-1 (2018)은 1억 1,700만 파라미터였지만, GPT-2 (2019)는 15억, GPT-3 (2020)는 1,750억 파라미터로 급성장했다. 이 과정에서 모델 크기, 데이터 크기, 연산량을 키우면 성능이 예측 가능하게 향상된다는 \keyterm{스케일링 법칙(scaling law)}이 발견되었다.

Kaplan et al. (2020)의 스케일링 법칙은 손실이 모델 크기 $N$, 데이터 크기 $D$, 연산량 $C$에 대해 멱법칙(power law)을 따른다는 것을 보였다:

\[
\mathcal{L}(N, D, C) \approx A \cdot N^{-\alpha_N} + B \cdot D^{-\alpha_D} + E \cdot C^{-\alpha_C} + L_\infty
\]

이 법칙의 실용적 시사점은 명확하다: 모델을 키우려면 데이터도 함께 키워야 한다. 100억 파라미터 모델을 10억 토큰으로 학습하면 과적합된다. Hoffmann et al. (2022)의 Chinchilla 논문은 파라미터 $N$당 약 20개의 토큰이 이상적이라고 제안했다.

\subsection{정렬의 등장 (2022$\sim$현재)}

GPT-3 (2020)는 1,750억 파라미터로 강력했지만, 사용자 질문에 도움이 되는 답변을 하지는 않았다. 그냥 ``다음 토큰''을 예측했을 뿐이다. 2022년 ChatGPT가 혁신적이었던 이유는 모델 크기가 아니라 \keyterm{정렬(alignment)}에 있었다.

인간 피드백을 활용한 강화학습(RLHF)이 GPT-3를 ``쓸모있는 챗봇''으로 변환시킨 것이다. 이후 Direct Preference Optimization (DPO), Constitutional AI 등 더 단순하고 효율적인 정렬 기법이 등장했다. 2권 8장에서 이 기법들을 직접 구현한다.

\section{LLM의 주요 응용}

LLM은 다양한 분야에서 활용되고 있다:

\subsection{자연어 처리 전통 태스크}
\begin{itemize}
  \item \textbf{기계 번역}: Google Translate, DeepL 등이 트랜스포머 기반.
  \item \textbf{요약}: 긴 문서를 짧게 요약. 뉴스, 논문, 회의록 등.
  \item \textbf{감성 분석}: 텍스트의 긍정/부정 판별. 리뷰, 소셜미디어 분석.
  \item \textbf{개체명 인식}: 텍스트에서 사람, 장소, 조직 등 추출.
\end{itemize}

\subsection{생성 태스크}
\begin{itemize}
  \item \textbf{대화}: ChatGPT, Claude, Gemini 등.
  \item \textbf{코드 작성}: GitHub Copilot, Cursor. Python, JavaScript 등 다양한 언어.
  \item \textbf{창작}: 소설, 시, 시나리오. AI 작가 도구.
  \item \textbf{이메일/문서 작성}: 비즈니스 이메일, 보고서 초안.
\end{itemize}

\subsection{전문 분야}
\begin{itemize}
  \item \textbf{의료}: 진단 보조, 의료 기록 요약, 환자 상담.
  \item \textbf{법률}: 판례 검색, 계약서 검토, 법률 자문.
  \item \textbf{교육}: 개인화 튜터, 문제 생성, 과제 채점.
  \item \textbf{금융}: 시장 분석, 리포트 작성, 사기 탐지.
\end{itemize}

\subsection{멀티모달}
최신 LLM은 텍스트뿐 아니라 이미지, 오디오, 비디오도 처리한다. GPT-4V, Gemini, Claude 3 등은 이미지를 이해하고 설명할 수 있다.

\section{왜 직접 만들어야 하는가}

``HuggingFace에서 모델을 다운로드하면 되는데, 왜 굳이 직접 만들어야 하는가?''라는 질문은 정당하다. 답은 세 가지다.

\textbf{첫째, 이해의 깊이}가 다르다. \code{model.generate()}를 호출하는 것과, 그 안에서 무슨 일이 일어나는지 아는 것은 차원이 다른 이해다. 직접 만들면 디버깅, 최적화, 커스터마이징이 모두 가능해진다. 추론 속도가 느릴 때 어디가 병목인지 진단할 수 있고, 새로운 아키텍처를 실험할 때 어디를 수정해야 할지 안다.

\textbf{둘째, 제어력}이 생긴다. 추론 속도를 높이고 싶을 때, 메모리를 줄이고 싶을 때, 새로운 아키텍처를 실험하고 싶을 때, 내부를 아는 사람과 모르는 사람의 대응력은 하늘과 땅 차이다. HuggingFace 모델을 그대로 쓰다가 이상한 출력이 나오면? 내부를 아는 사람은 어디를 디버깅해야 할지 안다.

\textbf{셋째, 미래 대비}다. LLM 분야는 빠르게 변한다. 새로운 논문이 매주 쏟아지는데, 그것을 이해하려면 기본기가 탄탄해야 한다. RoPE, GQA, SwiGLU 같은 최신 기법도 결국 어텐션, FFN, 정규화의 변형이다. 이 책을 통해 기본기를 다지면, 2권 \textit{Make LLM-advanced}에서 다루는 분산 학습, 양자화, 정렬 같은 고급 주제도 자연스럽게 흡수할 수 있다.

\begin{warnbox}[이 책의 한계]
이 책은 ``실제 ChatGPT를 만드는'' 것을 목표로 하지 않는다. 1권에서 만드는 모델은 1천만$\sim$1억 파라미터 수준의 소규모 GPT로, 품질은 GPT-2 small 정도다. 하지만 그 구조는 GPT-3, GPT-4와 본질적으로 같다. ``크기''는 2권에서 다루고, ``구조와 원리''를 1권에서 다룬다.

왜 큰 모델을 만들지 않는가? 175B 모델을 학습하려면 수백 개 GPU가 수주일 필요하고, 수억 달러 비용이 든다. 이것은 책으로 따라하기 어렵다. 대신 작은 모델로 원리를 익히면, 큰 모델도 같은 원리라는 것을 알게 된다.
\end{warnbox}

\section{이 책이 만들 모델}

1권을 완독하면 독자는 다음과 같은 모델을 직접 만들게 된다.

\begin{itemize}
  \item \textbf{아키텍처}: GPT 스타일 트랜스포머 디코더 (6$\sim$12층, 8헤드, 512$\sim$768 임베딩 차원)
  \item \textbf{토크나이저}: BPE (Byte Pair Encoding)와 Unigram 둘 다 직접 구현
  \item \textbf{학습}: AdamW + cosine LR 스케줄 + warmup
  \item \textbf{생성}: Greedy, Top-K, Top-P, Temperature 샘플링 전부 구현
  \item \textbf{평가}: Cross-entropy loss와 Perplexity 계산
\end{itemize}

2권 \textit{Make LLM-advanced}에서는 이 기반 위에 분산 학습(DDP, FSDP, ZeRO), 최신 아키텍처(RoPE, GQA, SwiGLU, MoE), 양자화(INT8/INT4/AWQ), 정렬(SFT, RLHF, DPO, LoRA), 추론 최적화(KV cache, PagedAttention)를 추가한다.

\section{학습 로드맵}

\begin{figure}[htbp]
\centering
\begin{tikzpicture}[
  every node/.style={font=\small\sffamily},
  box/.style={draw=inkblue, fill=inkblue!8, rounded corners=2pt, minimum width=3cm, minimum height=0.8cm, align=center},
  arr/.style={-{Stealth[length=4pt]}, inkblue!70, thick}
]
  \node[box] (ch1) at (0, 0) {1장\\도입};
  \node[box] (ch2) at (3.5, 0) {2장\\환경 세팅};
  \node[box] (ch3) at (7, 0) {3장\\토크나이저};
  \node[box] (ch4) at (10.5, 0) {4장\\임베딩};
  \node[box] (ch5) at (0, -1.8) {5장\\어텐션};
  \node[box] (ch6) at (3.5, -1.8) {6장\\트랜스포머};
  \node[box] (ch7) at (7, -1.8) {7장\\GPT 학습};
  \node[box] (ch8) at (10.5, -1.8) {8장\\생성};
  \node[box, fill=inkred!10, draw=inkred] (ch9) at (5.25, -3.6) {9장\\정리};

  \draw[arr] (ch1) -- (ch2);
  \draw[arr] (ch2) -- (ch3);
  \draw[arr] (ch3) -- (ch4);
  \draw[arr] (ch4) -- (ch5);
  \draw[arr] (ch5) -- (ch6);
  \draw[arr] (ch6) -- (ch7);
  \draw[arr] (ch7) -- (ch8);
  \draw[arr] (ch8) -- (ch9);
\end{tikzpicture}
\caption{1권 학습 로드맵. 각 장은 이전 장의 모듈을 import하여 사용.}
\label{fig:learning-roadmap}
\end{figure}

\section{요약}

\begin{itemize}
  \item 언어 모델은 텍스트 시퀀스에 확률을 할당하는 모델이며, 다음 토큰 예측이 핵심 목표다.
  \item N-gram에서 RNN/LSTM을 거쳐 2017년 트랜스포머가 등장하며 LLM 시대가 열렸다.
  \item 스케일링 법칙에 따라 모델, 데이터, 연산을 키우면 성능이 예측 가능하게 향상한다.
  \item 2022년 이후 정렬(RLHF/DPO)이 ``단순 LM''을 ``쓸모있는 챗봇''으로 변환시켰다.
  \item 직접 만들면 이해의 깊이, 제어력, 미래 대비 능력이 모두 확보된다.
  \item 1권에서는 소규모 GPT를, 2권에서는 ChatGPT 수준의 기법을 다룬다.
\end{itemize}

\section{연습 문제}

\begin{enumerate}
  \item ChatGPT, Claude, Gemini 중 하나를 사용해 다음 질문에 답해보라. ``언어 모델이 다음 토큰을 예측하도록 학습된다면, 어떻게 질문에 답할 수 있는가?'' 모델의 답변을 평가해보라.
  \item 다음 문장의 다음 단어 확률을 사람이라면 어떻게 매길지 생각해보라: ``태양은 동쪽에서 ...''. ``솟는''이 99\% 이상일 것이다. 왜 이 확률이 높은가?
  \item N-gram 모델의 한계 중 ``희소성 문제''를 설명하라. 학습하지 않은 3-gram ``북극곰이 춤을'' 다음 단어를 N-gram은 어떻게 예측할까?
  \item 트랜스포머가 RNN보다 병렬 처리에 유리한 이유를 설명하라.
  \item 스케일링 법칙에 따르면 모델을 10배 키우면 데이터는 얼마나 키워야 하는가? (Chinchilla 기준)
\end{enumerate}

다음 장에서는 본격적인 구현을 위해 Python, PyTorch, 그리고 이 책의 코드 저장소를 세팅한다.
